\documentclass[11pt, letterpaper]{article}
\usepackage[margin=0.9in]{geometry}
\usepackage{hyperref}
\usepackage{enumitem}
\usepackage{xcolor}
\usepackage{titlesec}
\usepackage{fancyhdr}
\usepackage{tcolorbox}
\usepackage{booktabs}

\definecolor{accent}{HTML}{1A5490}
\definecolor{qbg}{HTML}{E8F0FE}
\hypersetup{colorlinks=true, linkcolor=accent, urlcolor=accent}

\pagestyle{fancy}
\fancyhf{}
\rhead{STA 561D | Duke}
\lhead{Infinite Connections --- FAQ}
\rfoot{\thepage}
\renewcommand{\headrulewidth}{0.4pt}

\titleformat{\section}{\large\bfseries\color{accent}}{\thesection.}{0.5em}{}
\titlespacing\section{0pt}{8pt}{4pt}

\newcommand{\qitem}[1]{\vspace{0.6em}\noindent\textbf{\color{accent}Q. #1}\\[0.2em]}
\newcommand{\aitem}{\noindent\textbf{A.} }

\title{\textbf{Infinite Connections}\\[0.2em]
\large Frequently Asked Questions}
\author{Abuzar \quad Adreama \quad Burak \quad Hengkai \quad Kaiwen\\
STA 561D --- Duke University}
\date{April 2026}

\begin{document}
\maketitle
\thispagestyle{fancy}
\vspace{-2em}

\section*{Introduction}
This FAQ anticipates the questions a careful, skeptical reader would ask after
reading our executive summary or technical report. We group them by topic: how
each of the five teammates' methods works, how we prevent duplicates, how we
measure quality, cost and scale, and known limitations.

% =====================================================================
\section*{The five methods, explained individually}

\qitem{What method did Kaiwen build?}
\aitem
Kaiwen implemented two pipelines head-to-head. \textbf{Pipeline A} is a
reference implementation of the iterative generator from the ``Making New
Connections'' paper --- five AI calls per puzzle, one for each of the four
groups plus a final editor pass. \textbf{Pipeline B} is a new design called
\textbf{Category-First Retrieval (CFR)}: instead of asking the AI to produce
words, it asks the AI only for \emph{category names}, and then a word-lookup
against a 14{,}877-word dictionary finds the best matching words. Mode~A of
CFR uses zero AI calls (it samples category names from the existing NYT
catalog); Mode~B uses one AI call per puzzle to invent fresh category names.
One AI call replaces five at the same quality.

\qitem{What method did Adreama build?}
\aitem
Adreama's method is a single-call iterative generator. A GPT-4.1 call returns
all four groups at once as structured JSON, with a long context prompt that
carries forward every word and category the generator has used across
sessions. The generator keeps three persistent memory files
(\texttt{used\_words.json}, \texttt{used\_categories.json},
\texttt{used\_boards.json}) so duplicates never creep in even over thousands
of puzzles. A library of 48 hand-curated ``category webs'' (tuples of
\textit{real} + \textit{trap} categories) seeds thematic false-group designs.
The generator retries up to four times on parse or validation errors; batches
of 50 attempts typically yield $\sim$35 valid puzzles.

\qitem{What method did Burak build?}
\aitem
Burak's method is almost entirely non-AI --- the AI only gets called to label
one group. The pipeline builds puzzles color by color, each color using a
different algorithm: \textbf{purple} is drawn from a pool scored by eight
rule-based mechanisms (e.g., collective nouns, contronyms), optionally with a
machine-learning classifier reranking the pool; \textbf{green} is built by
rhyme anchors using the CMU pronunciation dictionary; \textbf{blue} is a
niche-category generator where Claude Haiku validates the group; and
\textbf{yellow} uses word2vec semantic clusters. ``Impostor'' words are
planted deliberately to create false-group pressure across colors.

\qitem{What method did Hengkai build?}
\aitem
Hengkai is the yellow-difficulty specialist. The generator is fully
algorithmic --- zero AI calls --- and uses pre-computed GloVe word embeddings,
ConceptNet relations, and corpus frequency filters. Candidates are scored
within several relation types (taxonomic, synonymy, verb--noun association,
ConceptNet link) and the highest-scoring combinations become yellow groups.
It is the only teammate's generator that validates each candidate group
against the full 554-puzzle NYT history before acceptance, rejecting any
group that shares three or more words with a past NYT group. A checkpoint
system lets batch runs of 700+ puzzles resume after interruption.

\qitem{What method did Abuzar build?}
\aitem
Abuzar's method is an agent-oriented system. A single GPT-4.1 call produces
the full 4-group puzzle in JSON, but the prompt is heavily engineered with
banned-category lists, banned-purple-type counters, and a trap-word
mechanism where one purple word is specifically designed to impersonate a
simpler category. The pipeline enforces self-uniqueness through persistent
word, category, and board-signature sets; it also ships with a Flask web
server and a ``concept inspirations'' library for theme variety. It does
\emph{not} check against the 554 past NYT puzzles --- only against its own
generated history.

\qitem{Why build five different methods instead of picking the best one?}
\aitem
The Connections puzzle is deceptively hard. A single method tends to be great
at one category style and weak at others: semantic-similarity approaches nail
``SYNONYMS FOR X'' but stumble on wordplay like ``\_\_\_FISH'' or hidden
patterns like ``WORDS CONTAINING COLORS.'' By spreading the work across five
complementary methods we cover more of the design space than any one
approach could. It also matches the professor's suggestion that we
\emph{identify the different mechanisms in which words are grouped and build
separate generators for each.}

\qitem{Why use embeddings at all --- aren't LLMs just as good?}
\aitem
LLMs are good at creative naming but inconsistent at ranking. Asked to pick
``the best 4 of these 8 words,'' an LLM gives different answers each time
depending on temperature, seed, and context. Embeddings (numerical
fingerprints of each word's meaning) give us \emph{deterministic} similarity
scores: we get the same answer every time, which makes difficulty colors
and validation reproducible. The combined approach --- LLM for naming,
embeddings for ranking --- gets the best of both.

% =====================================================================
\section*{Preventing duplicates}

\qitem{How does each teammate prevent duplicates against past NYT puzzles?}
\aitem
Not all five methods check against the 554 past NYT puzzles. Here's the
honest breakdown:

\begin{center}
\small
\begin{tabular}{p{0.15\linewidth} p{0.5\linewidth} p{0.3\linewidth}}
\toprule
\textbf{Teammate} & \textbf{Check vs.~past NYT puzzles?} & \textbf{Self-duplicate check} \\
\midrule
Kaiwen  & \textbf{Yes.} 16-word set overlap $>\!6$ flag + exact 4-word group match against all $\sim$2{,}216 past NYT groups. Enforced during generation (CFR skips any candidate that matches a past group). & Implicit via per-batch unique words \\
Hengkai & \textbf{Yes.} Rejects any yellow group that shares 3 or more words with a past NYT group. & Anchor-repeat limit + checkpointing \\
Adreama & No explicit check against NYT puzzles. & \textbf{Yes:} persistent JSON files \\
Abuzar  & No explicit check against NYT puzzles. & \textbf{Yes:} persistent JSON files \\
Burak   & No explicit check against NYT puzzles. & In-memory only (per notebook run) \\
\bottomrule
\end{tabular}
\end{center}

The strongest guarantee comes from Kaiwen's pipeline, which enforces the check
both during word selection (skip-and-retry) and as a final gate. Because
Kaiwen's roundtable validator is used as the shared acceptance step across
the team's final merged dataset, any puzzle that slips through a teammate's
own generator still has to pass the 4-word-group match before it lands in the
submitted dataset. In that sense the team-wide submitted dataset \emph{is}
guarded against NYT duplication --- via Kaiwen's step --- even when the
upstream generator didn't check.

\qitem{What about self-duplication across batches?}
\aitem
Four of the five methods have an explicit self-duplication defense. Adreama
and Abuzar persist three files each (\texttt{used\_words.json},
\texttt{used\_categories.json}, \texttt{used\_boards.json}) that carry state
across sessions. Hengkai's batch runner checkpoints to a JSON file and
resumes without repeating groups. Kaiwen's CFR excludes already-used words
within a puzzle and deduplicates by board signature across batches. Burak's
notebook, the exception, holds its dedup state only in memory --- a repeated
notebook run can in principle regenerate the same puzzle. In practice our
merged dataset is deduplicated by board signature as a final step, which
removes any cross-teammate collisions.

\qitem{What about near-duplicates --- a puzzle with 5 overlapping words?}
\aitem
Kaiwen's threshold is 6 (any overlap of 7 or more words flags the puzzle;
Hengkai's threshold for yellow groups alone is 3). With 16-word puzzles
drawn from a 14{,}877-word bank, overlapping 5 or 6 words happens by chance
perhaps once in thousands of generated puzzles. The per-puzzle metadata we
save includes the overlap count against each past NYT puzzle so a reviewer
can audit this directly and adjust the threshold if desired.

\qitem{Could the LLM leak a memorized NYT puzzle?}
\aitem
The underlying language models were trained on web text, so they have
certainly seen Connections discussions. Three of our pipelines
(Kaiwen, Burak, Hengkai) do not let the LLM produce the final word list
directly --- they only let it propose category names or bless a candidate
group, while the actual words come from deterministic dictionary lookup. For
the pipelines that do ask the LLM to emit the full 16 words (Adreama and
Abuzar), the group-level dedup in our merged pipeline catches any leaked
group before acceptance.

% =====================================================================
\section*{Quality and evaluation}

\qitem{99\% pass rate --- doesn't that mean your solver is too lenient?}
\aitem
We don't trust a single solver; that's why Kaiwen's pipeline runs \emph{two}
independent ones. A puzzle is accepted only if at least one solver fully
recovers the intended 4-group partition. Both solvers use deterministic
algorithms (greedy cosine selection and beam-search clustering) and solve
only 2--3\% of the real NYT puzzles when tested directly on them --- they
are conservative, not easy. So a puzzle that our solvers \emph{can} crack
has a cleanly recoverable structure (exactly what a good puzzle should have).

\qitem{How close to the NYT style are the generated puzzles?}
\aitem
We have not yet run a formal human evaluation. Internally, on the categories
where our methods excel (synonyms, types-of-X, themed nouns) many puzzles
are difficult to distinguish from NYT output on casual inspection. On pure
wordplay categories (``\_\_\_FISH,'' letter homophones), the output is
weaker --- see the limitations section. We added a thumbs-up/thumbs-down
rating button to the web app so real players can rate puzzles at scale, and
we plan to collect at least several hundred ratings before submission.

\qitem{What's the unique-solution guarantee?}
\aitem
A Connections puzzle is only valid if exactly one 4-way partition fits the
intended themes. Our pipeline enforces this in two ways: (1) during
generation, CFR's candidate selector skips any 4-word group with a high
cross-group ``bleed'' score (words that fit another group too well), and (2)
during validation, our two solvers must find the intended partition --- if
they find a \emph{different} partition that also satisfies the themes, the
puzzle is rejected as ambiguous.

\qitem{Is the word bank big enough?}
\aitem
The NYT has used about 4{,}918 unique words across all 554 of its puzzles.
Kaiwen's augmented bank has \textbf{14{,}877 words} --- three times larger.
In the CFR output, 62--69\% of generated words come from the \emph{non-NYT}
portion, showing the expansion is actually being used. Burak and Hengkai
additionally pull from WordNet, ConceptNet, and the Brown corpus, which
expand the vocabulary further within those specialized pipelines.

% =====================================================================
\section*{Cost and scale}

\qitem{What does it cost to generate 10{,}000 puzzles with each method?}
\aitem
\begin{center}
\begin{tabular}{lr}
\toprule
\textbf{Method} & \textbf{Approx.~cost / 10{,}000 puzzles} \\
\midrule
Kaiwen CFR Mode A (no AI calls)        & \$0      \\
Kaiwen CFR Mode B (1 AI call each)     & \$3      \\
Hengkai yellow generator (no AI calls) & \$0      \\
Burak pattern-builder (2 AI calls/puzzle) & ${\sim}\$20$  \\
Adreama iterative (1 long AI call/puzzle) & ${\sim}\$10$  \\
Abuzar agent system (1 long AI call/puzzle) & ${\sim}\$15$ \\
Pipeline A baseline (5 AI calls/puzzle)  & ${\sim}\$50$ \\
\bottomrule
\end{tabular}
\end{center}
The cheapest methods let us generate unlimited puzzles at zero cost; even
the most expensive is within a routine student budget.

\qitem{Why are different methods cheaper than others?}
\aitem
They make fewer calls to the paid AI service. Pipeline A asks the AI five
times per puzzle. Adreama and Abuzar each make one long call. CFR Mode B
makes one short call. CFR Mode A and Hengkai's generator make zero calls ---
they use dictionaries, embeddings, and math, which run free on a laptop.

\qitem{Will this scale to 100{,}000 puzzles?}
\aitem
Yes. CFR Mode A runs at about 1.3 seconds per puzzle on a single laptop CPU
with no network calls, so 100{,}000 puzzles take about 36 hours of wall time
and cost \$0. Hengkai's batch generator scales similarly. The bottleneck
at that scale is the roundtable validator, not generation.

% =====================================================================
\section*{Reproducibility}

\qitem{Can someone else reproduce your results without API keys?}
\aitem
Mostly, yes. Kaiwen's code ships with a \texttt{DRY\_RUN=true} default:
every call to the AI service is intercepted and replaced by a realistic
mock response. The master notebook, the web app, the test suite, and the
solver benchmark all run end-to-end with no keys configured. Setting
\texttt{DRY\_RUN=false} (with \texttt{OPENAI\_API\_KEY}) switches to real
calls. Adreama's and Abuzar's generators require live API access. Burak's
and Hengkai's generators run offline after pre-computing their embeddings
and corpora.

\qitem{How do I re-run the benchmarks?}
\aitem
Two commands for Kaiwen's pipeline:
\begin{verbatim}
python scripts/cfr/generate_cfr.py --mode remix --count 100
python scripts/cfr/generate_cfr.py --mode fresh --count 100
\end{verbatim}
Output is written to \texttt{data/generated/cfr/\dots} with per-puzzle
metadata (pass rate, solver traces, dedup results) for audit. Other
teammates' generators live in their own folders and each ship with
instructions.

\qitem{Where is the code?}
\aitem
The primary repository is at
\url{https://github.com/Kevin2330/nyt-connections-generator}. The live web
app at
\url{https://kevin2330.github.io/nyt-connections-generator/}
loads 540+ validated puzzles from the same repo. Individual teammate code
lives in the \texttt{STA561/} directory of the submission; Abuzar's code
additionally lives in a private team repository.

% =====================================================================
\section*{Limitations}

\qitem{What's the biggest weakness of the overall system?}
\aitem
Wordplay and syntactic categories. Our embedding-based methods (Kaiwen,
Hengkai) rely on semantic similarity, which captures ``SYNONYMS FOR SMART''
but not ``WORDS THAT END IN -ING'' or ``\_\_\_FISH.'' These pattern-based
categories require string-level and phonetic reasoning that embeddings
don't naturally do well. Burak's CMU-rhyme green generator and Abuzar's
multi-agent system partly address this, but the hardest (purple) categories
remain our collective weakest output.

\qitem{Why does WordNet sometimes return awkward words (e.g., MORPHOPHONEMIC)?}
\aitem
WordNet includes many archaic and technical words. Kaiwen's pipeline filters
by corpus frequency and character length, but some oddities slip through. A
higher frequency threshold or intersection with a common-English list (the
Google 10{,}000-most-common-words list is already included as an optional
filter) would tighten this further.

\qitem{Why does Burak's generator have no cross-session deduplication?}
\aitem
It is an honest gap. Burak's notebook holds its used-words set only in
memory, so two separate notebook runs could in principle generate the same
puzzle. In practice the final merged dataset is deduplicated by
16-word board signature as a post-processing step, which removes any such
collisions before submission.

% =====================================================================
\section*{Teamwork}

\qitem{Who did what?}
\aitem
\begin{center}
\begin{tabular}{ll}
\toprule
\textbf{Teammate} & \textbf{Contribution} \\
\midrule
Adreama & Iterative AI generator with cross-session memory and trap-word webs \\
Burak   & Non-AI pattern-driven builder (purple mechanism + rhyme green + word2vec yellow) \\
Hengkai & Yellow specialist: embedding + ConceptNet scoring, NYT-history overlap check \\
Kaiwen  & Pipeline A baseline, CFR Modes A \& B, roundtable validator, web app \\
Abuzar  & Multi-agent critic system + Flask web server + concept-inspiration library \\
\bottomrule
\end{tabular}
\end{center}

\end{document}
